\thispagestyle{plain}
\noindent{\small\sffamily\color{GearTeal}GEARPRO\quad /\quad 阶段研究报告}\par
\vspace{1.1cm}
\noindent{\fontsize{29}{37}\selectfont\sffamily\bfseries\color{GearInk}齿轮外观缺陷检测\par}
\vspace{0.2cm}
\noindent{\fontsize{23}{31}\selectfont\sffamily\bfseries\color{GearInk}系统技术报告\par}
\vspace{0.55cm}
\noindent{\large\color{GearMuted}自建数据 · 两类缺陷选型 · 参数依据 · 软件集成}\par
\vspace{0.4cm}
\noindent{\small\color{GearMuted}GearPro\quad 阶段记录日期：2026年10月7日}\par
\vspace{0.55cm}
\noindent\textcolor{GearTeal}{\rule{\linewidth}{1.1pt}}\par
\vspace{0.4cm}
\noindent{\large\sffamily\bfseries 摘要}\par\medskip
本项目面向齿轮外观检查，重点检测划痕和局部齿部缺损。我们自行采购或定制50个齿轮，拍摄多角度图像并标注可见缺陷。在早期分类实验的基础上，分别比较两类缺陷的标签划分、检测结构和模型组合，利用验证集选择模型与判定参数。最终方案将分类模型与目标检测模型的分数加权组合，再合并两类缺陷的判断。在150张离线测试图像上，检出63/71张缺陷图像，召回率为88.73\%；误报18/79张无可见目标缺陷图像，误报率为22.78\%。相比已测试的统一缺陷分类模型，检出数净增加9张，误报数减少15张；两套方案的训练数据和计算预算不同。项目已完成软件集成和第一代原型联调，目前使用上下两台工业相机和Jetson NX。新模型只做过简单现场测试，没有系统记录的整机性能数据。现阶段面向辅助人工复核，后续还需用新齿轮和现场测试检验效果。

\textbf{关键词：} 齿轮外观检测；划痕；齿部缺损；模型融合；边缘计算

\ReportNote{
阅读说明：正文围绕任务、方法、实验依据和结果展开；完整配置与补充分析放在附录。项目事实、实验记录与历史报告的来源说明见\protect\hyperlink{appendix-b}{附录B}；学术文献采用数字编号引用。本报告未新增训练或现场测试。
}
